\centerline{\bf \Huge Графы F}

\problem{\textit{\textbf{1}}} Найдите сумму степеней вершин изображенного на рисунке графа и уменьшите найденную сумму на количество ребер графа.

\problem{\textit{\textbf{2}}} У графа 15 вершин степени 7 и еще 11 вершин степени 13. Других вершин в этом графе нет. Сколько ребер в этом графе?

\problem{\textit{\textbf{3}}} В графе 19 вершин, каждая из которых имеет индекс 6. Других вершин в этом графе нет. Сколько у него ребер?

\problem{\textit{\textbf{4}}} В графе 123 ребра, а каждая вершина имеет индекс 3. Других вершин в этом графе нет. Сколько у него вершин?

\problem{\textit{\textbf{5}}} Какие графы, изображенные на рисунке, можно нарисовать, не отрывая карандаша от бумаги и проводя каждое ребро ровно один раз?

\problem{\textit{\textbf{6}}} Какой наименьшей длины должна быть проволока, чтобы из неё можно было сложить рёберную модель куба с ребром 23 см? Ответ укажите в метрах.

\problem{\textit{\textbf{7}}} Какое наименьшее число рёбер придется пройти дважды, чтобы обойти все рёбра тетраэдра и вернуться в исходную вершину?

\problem{\textit{\textbf{8}}} Граф, не содержащий ни одной замкнутой ломаной, называется лесом. Пусть лес состоит из 67 деревьев и имеет А вершин и Б ребер. Чему равно А - Б?

\problem{\textit{\textbf{9}}} В графе 132 ребра. Каждая вершина графа имеет или степень 8, или степень 14. Причём вершин степени 8 столько же, сколько вершин степени 14. Сколько вершин в этом графе?

\problem{\textit{\textbf{10}}} В графе 10 рёбер. Каждая вершина графа имеет или степень 2, или степень 3. Причём вершин степени 2 на 5 меньше, чем вершин степени 3. Сколько вершин в этом графе?

\problem{\textit{\textbf{11}}} Оля нарисовала схему, не отрывая карандаша от листа бумаги и не проводя никакую линию дважды. В какой точке Оля закончила рисовать схему, если она начала её рисовать в точке E?

\problem{\textit{\textbf{12}}} В городе Маленьком 15 телефонов. Можно ли их соединить проводами так, чтобы каждый телефон был соединен ровно с пятью другими?

\newpage
\hspace{3mm}

\problem{\textit{\textbf{13}}} В Тридевятом царстве лишь один вид транспорта — ковер-самолет. Из столицы выходит 21 ковролиния, из города Дальний — одна, а из всех остальных городов — по 20. Можно ли из столицы долететь в Дальний (возможно, с пересадками).

\problem{\textit{\textbf{14}}} В стране Семерка 15 городов, каждый из которых соединен дорогами не менее, чем с семью другими. Верно ли, что из любого города можно ли добраться до любого другого, возможно, проезжая через другие города?